\chapter{From study to research}
\label{ch:research}

\section{Reading the project's own results}

The pilot is the best exercise in this guide, because its results are mixed and honest. With matched heads,
\method{} trails the simple concatenation baseline by 0.3 to 1.5 AP on clean LLVIP images, close to run-to-run
noise. It holds up better when the thermal image is shifted, and when told that the thermal camera is gone it
still detects, where concatenation collapses. When the thermal camera is switched off without warning, both
models find no one, although the people are plainly visible. Each of these is a research question, and each
can be stated as a hypothesis, tested with one controlled run, and decided by a rule written down in advance.

\section{Two lessons from the training logs}

The first full training run diverged to NaN from its second epoch. The cause was half precision: over twelve
thousand tokens, the scan's inputs and outputs left float16's range. The fix was to keep the whole state-space
path in float32. A later run diverged again once modality dropout forced one camera's weight to zero and the
model pushed the other towards zero as well: the fusion divided by almost nothing, the gradient became a million
times too large, and float16 overflowed. The fix was to compute the weights in float32 and floor their total.
Both bugs were invisible in the unit tests and obvious in the logs. Read the logs.

\section{A weekly rhythm}

\begin{tabularx}{\textwidth}{@{}P{3cm}L@{}}
\hd{day} & \hd{activity} \\
Monday & one lecture and its notes \\
Tuesday & the matching exercise or assignment \\
Wednesday & one paper, read three times \\
Thursday & reproduce one figure or number from that paper, however small \\
Friday & read the matching code in this repository and change one thing \\
Weekend & write a one-page summary of the week in your own words \\
\end{tabularx}

\section{When you are stuck}

Shrink the problem until it is small enough to see: one image instead of a dataset, one block instead of a
network, ten steps instead of thirty epochs. Make a model overfit a single batch before training it properly.
Change one thing at a time. When a result surprises you, suspect the code before the theory, and the data before
the code.
